\documentclass[11pt]{article}
\usepackage[letterpaper,margin=0.9in]{geometry}
\usepackage{amsmath,amssymb,amsthm}
\usepackage{titlesec}
\usepackage[hidelinks]{hyperref}
\titleformat{\section}{\normalsize\bfseries}{\thesection.}{0.5em}{}
\titlespacing*{\section}{0pt}{9pt}{3pt}
\setlength{\parskip}{4pt}
\setlength{\parindent}{0pt}

\newtheorem{theorem}{Theorem}
\newtheorem{lemma}[theorem]{Lemma}
\newcommand{\dist}{\delta}
\newcommand{\R}{\mathbb{R}}

\title{\vspace{-2.2em}\textbf{Rollin' Down the Biggest Hill: Analysis Draft}\\
\large Correctness and running time of Dijkstra's algorithm\\
\normalsize CS 351 Analysis of Algorithms, Project Checkpoint}
\author{The \textsc{wheel} Ones: Conor Vahey, Felicity Chin, Oscar Yaun}
\date{\vspace{-1.5em}}

\begin{document}
\maketitle

\section{Model}

A world is a directed graph $G=(V,E)$. Our test worlds are $n\times n$ grids: each cell is a node with an elevation $\mathrm{el}(v)\in\R$, and each cell has an edge in each direction to each of its four neighbours. The cost of an edge is
\[
w(u,v)=\begin{cases}0 & \text{if } \mathrm{el}(v)<\mathrm{el}(u)\ \ \text{(downhill: roll)}\\ 1 & \text{otherwise}\ \ \text{(flat or uphill: push)}.\end{cases}
\]
The cost of a path is its number of pushes. $\dist(s,t)$ is the minimum cost of any path from $s$ to $t$. Every edge has a reverse edge, so every node can reach every other node. On a grid $|E| = 4n(n-1) < 4|V|$. All costs are non-negative, which is the property Dijkstra's algorithm needs.

We use the \emph{subpath property}: if $P$ is a shortest path from $s$ to $u$ and $y$ lies on $P$, then the part of $P$ from $s$ to $y$ is a shortest path from $s$ to $y$, since otherwise swapping in a cheaper path would make $P$ cheaper \cite[Lemma~22.1]{clrs}.

\section{The algorithm as implemented}

(\texttt{heely/dijkstra.py}.) Each node $v$ has a tentative distance $d[v]$, initially $\infty$ except $d[s]=0$, and a flag \emph{settled}, initially false. A binary min-heap holds pairs $(d, v)$. Repeat until the heap is empty: pop the pair with the smallest $d$; if $v$ is already settled, skip it; otherwise mark $v$ settled, and stop if $v=t$; otherwise for each edge $(v,x)$, if $d[v]+w(v,x)<d[x]$, set $d[x]\gets d[v]+w(v,x)$, set $\mathrm{parent}[x]\gets v$, and push $(d[x],x)$. Instead of decreasing a key in the heap we push a new pair and skip old ones when they come out (\emph{lazy deletion}), which is the standard way to use Python's \texttt{heapq}. Otherwise this is the algorithm in CLRS Section~22.3 \cite{clrs}.

\section{Correctness}

\begin{lemma}\label{lem:lazy}
When a pair $(d,v)$ is popped and $v$ is not yet settled, $d=d[v]$, and $d[v]\le d[x]$ for every unsettled node $x$.
\end{lemma}
\begin{proof}
Let $(d,v)$ be an arbitrary popped pair with $v$ unsettled. Every time $d[v]$ was lowered, a pair with the new value was pushed, so the heap holds $(d[v],v)$, and $d[v]\le d$ because $d[v]$ only decreases. Since the heap pops the smallest key and $(d,v)$ came out, $d\le d[v]$, so $d=d[v]$. In the same way every unsettled $x$ with $d[x]<\infty$ has a pair $(d[x],x)$ in the heap, and the heap returned the smallest, so $d[v]\le d[x]$. Because $(d,v)$ was arbitrary, the lemma holds for every pop.
\end{proof}

\begin{theorem}[CLRS Theorem 22.6]\label{thm:dij}
When a node $u$ is settled, $d[u]=\dist(s,u)$. In particular, when the algorithm stops at $t$, $d[t]=\dist(s,t)$, and following parent pointers from $t$ back to $s$ gives a path of cost $\dist(s,t)$.
\end{theorem}
\begin{proof}
First, $d[u]\ge\dist(s,u)$ always, because $d[u]$ is either $\infty$ or the cost of a real path built by relaxations, and no path costs less than the distance. Now suppose, for contradiction, that some node is settled with $d$ larger than its distance, and let $u$ be the first such node. Then $u\neq s$, since $d[s]=0=\dist(s,s)$. Let $P$ be a shortest path from $s$ to $u$. When $u$ is settled, $s$ is settled and $u$ is not, so let $y$ be the first node on $P$ that is not settled, and $x$ the node just before it, which is settled. Since $u$ is the first bad node, $d[x]=\dist(s,x)$. When $x$ was settled the edge $(x,y)$ was relaxed, so $d[y]\le\dist(s,x)+w(x,y)=\dist(s,y)$ by the subpath property, and together with $d[y]\ge\dist(s,y)$ this gives $d[y]=\dist(s,y)$. Since $y$ is on $P$ no later than $u$ and costs are non-negative, $\dist(s,y)\le\dist(s,u)$. So $d[y]\le\dist(s,u)<d[u]$. But $u$ was popped while $y$ was unsettled, so by Lemma~\ref{lem:lazy} $d[u]\le d[y]$, a contradiction. So no such $u$ exists.

For the path: when $\mathrm{parent}[x]$ is set to $v$, $v$ is settled, so $d[v]$ is final and $d[x]=d[v]+w(v,x)$. Once $x$ is settled, $d[x]$ and $\mathrm{parent}[x]$ never change. So walking parents back from $t$ passes through settled nodes, each settled strictly earlier than the one before, so the walk cannot loop, and it can only end at $s$, the one settled node without a parent. Each step lowers $d$ by exactly the edge cost, so the path costs $d[t]-d[s]=\dist(s,t)$.
\end{proof}

\section{Running time}

\begin{theorem}\label{thm:time}
Dijkstra's algorithm with a binary heap and lazy deletion runs in $O((|V|+|E|)\log|V|)$ time and $O(|V|+|E|)$ space. On an $n\times n$ grid this is $O(n^2\log n)$.
\end{theorem}
\begin{proof}
Each node is settled at most once, and a node's outgoing edges are relaxed only when it is settled, so there are at most $|E|$ relaxations. That means at most $|E|+1$ pushes and at most $|E|+1$ pops. The heap never holds more than $|E|+1$ pairs, so each push or pop costs $O(\log|E|)=O(\log|V|)$, since $|E|\le|V|^2$. Setting up the arrays costs $O(|V|)$. The total is $O(|V|+|E|\log|V|)=O((|V|+|E|)\log|V|)$. The arrays take $O(|V|)$ space, the adjacency lists $O(|V|+|E|)$, and the heap $O(|E|)$. On a grid $|E|<4|V|=4n^2$, so the time is $O(n^2\log n)$.
\end{proof}

Stopping when $t$ is settled does not change the worst case, but in practice it means only the nodes closer to $s$ than $t$ (plus some ties) are settled. The preliminary results show about half of all nodes being settled on average, which is what the second algorithm, ALT, aims to reduce.

\section*{Sources}
\begin{thebibliography}{9}
\bibitem{clrs} Cormen, T. H., Leiserson, C. E., Rivest, R. L., \& Stein, C. (2022). \emph{Introduction to Algorithms} (4th ed.). MIT Press. Chapter 22, Lemma 22.1 and Theorem 22.6.
\end{thebibliography}

\end{document}
